% ======================================================================
%  PART 4 -- Parametric Estimation: Yule--Walker and Least Squares
% ======================================================================
\section{Parametric Estimation: Yule--Walker and Least Squares}
\label{sec:estimation}

Having modelled $\{X_t\}$ as an ARMA process, we now \emph{fit} it: given data
$X_1,\dots,X_n$, estimate the autoregressive coefficients $\phi_1,\dots,\phi_p$,
the moving-average coefficients $\theta_1,\dots,\theta_q$ and the innovation
variance $\sigma_\eps^2$. This part develops the two workhorse estimators of the
course. The \textbf{Yule--Walker} (method-of-moments) estimator turns the AR
recursion into a linear system in the autocovariances and is solved by a single
matrix inversion (accelerated by the Levinson--Durbin algorithm). The
\textbf{least-squares} estimator minimises the sum of squared one-step
prediction errors; for pure AR models it is an ordinary linear regression
(forward, backward, or forward--backward), and for general ARMA models it
becomes a nonlinear minimisation in which the unobserved innovations
$\{\eps_t\}$ are reconstructed recursively. Throughout, $\eps_t\iid\mathcal
N(0,\sigma_\eps^2)$ is Gaussian white noise and the AR part is assumed causal
(roots of $\Phi(z)$ outside the unit circle), so that $\eps_t$ is uncorrelated
with the past $X_{t-k}$, $k>0$ --- the single fact that makes the whole
derivation work.

% ----------------------------------------------------------------------
\subsection{Definitions}

\begin{definition}{Yule--Walker estimators (AR$(p)$)}{p4-yw}
Let $\{X_t\}$ be a mean-zero causal AR$(p)$ process
$X_t=\sum_{k=1}^p\phi_kX_{t-k}+\eps_t$, observed at $X_1,\dots,X_n$. Form the
sample ACVS (mean known to be zero)
\[
\hat\acvs_\tau=\frac1n\sum_{t=1}^{n-\abs\tau}X_tX_{t+\abs\tau},
\]
collect $\hat{\boldsymbol\acvs}_p=(\hat\acvs_1,\dots,\hat\acvs_p)\Tr$ and the
symmetric Toeplitz matrix
\[
\hat\Gamma_p=
\begin{pmatrix}
\hat\acvs_0 & \hat\acvs_1 & \cdots & \hat\acvs_{p-1}\\
\hat\acvs_1 & \hat\acvs_0 & \cdots & \hat\acvs_{p-2}\\
\vdots & \vdots & \ddots & \vdots\\
\hat\acvs_{p-1} & \hat\acvs_{p-2} & \cdots & \hat\acvs_0
\end{pmatrix}.
\]
The \textbf{Yule--Walker estimators} are
\[
\hat{\boldsymbol\phi}_p=\hat\Gamma_p^{-1}\hat{\boldsymbol\acvs}_p,
\qquad
\hat\sigma_\eps^2=\hat\acvs_0-\sum_{j=1}^p\hat\phi_j\,\hat\acvs_j .
\]
They are the empirical analogue of the population identities
$\boldsymbol\acvs=\Gamma\boldsymbol\phi$ and
$\acvs_0=\sum_j\phi_j\acvs_j+\sigma_\eps^2$.
\end{definition}

\begin{definition}{Forward least-squares estimators (AR$(p)$)}{p4-fls}
Write the AR equations for the observable times $t=p+1,\dots,n$ as the linear
model $\mathbf X_F=F\boldsymbol\phi+\boldsymbol\eps$, where
\[
\mathbf X_F=\begin{pmatrix}X_{p+1}\\ \vdots\\ X_n\end{pmatrix},
\quad
F=\begin{pmatrix}
X_p & X_{p-1} & \cdots & X_1\\
X_{p+1} & X_p & \cdots & X_2\\
\vdots & \vdots & & \vdots\\
X_{n-1} & X_{n-2} & \cdots & X_{n-p}
\end{pmatrix},
\quad
\boldsymbol\phi=\begin{pmatrix}\phi_1\\ \vdots\\ \phi_p\end{pmatrix}.
\]
The \textbf{forward least-squares estimator} minimises
\[
\mathrm{SS}_F(\boldsymbol\phi)=\norm{\mathbf X_F-F\boldsymbol\phi}^2
=\sum_{t=p+1}^n\Big(X_t-\sum_{k=1}^p\phi_kX_{t-k}\Big)^2
=\sum_{t=p+1}^n\eps_t^2,
\]
and equals the ordinary least-squares solution
\[
\hat{\boldsymbol\phi}_F=(F\Tr F)^{-1}F\Tr\mathbf X_F,
\qquad
\hat\sigma_F^2=\frac{\norm{\mathbf X_F-F\hat{\boldsymbol\phi}_F}^2}{\,n-2p\,}.
\]
The first $p$ observations $X_1,\dots,X_p$ are discarded as conditioning values.
\end{definition}

\begin{definition}{Backward least-squares estimators (AR$(p)$)}{p4-bls}
Exploiting time reversibility of a Gaussian AR process, regress the
\emph{first} $n-p$ values on their \emph{future} lags: for $t=1,\dots,n-p$,
$X_t=\sum_{j=1}^p\phi_jX_{t+j}+\nu_t$. In matrix form
$\mathbf X_B=B\boldsymbol\phi+\boldsymbol\nu$ with
\[
\mathbf X_B=\begin{pmatrix}X_1\\ \vdots\\ X_{n-p}\end{pmatrix},
\qquad
B=\begin{pmatrix}
X_2 & X_3 & \cdots & X_{p+1}\\
X_3 & X_4 & \cdots & X_{p+2}\\
\vdots & \vdots & & \vdots\\
X_{n-p+1} & X_{n-p+2} & \cdots & X_n
\end{pmatrix}.
\]
The \textbf{backward least-squares estimator} minimises
$\mathrm{SS}_B(\boldsymbol\phi)=\norm{\mathbf X_B-B\boldsymbol\phi}^2
=\sum_{t=1}^{n-p}\big(X_t-\sum_{k=1}^p\phi_kX_{t+k}\big)^2$ and is
\[
\hat{\boldsymbol\phi}_B=(B\Tr B)^{-1}B\Tr\mathbf X_B,
\qquad
\hat\sigma_B^2=\frac{\norm{\mathbf X_B-B\hat{\boldsymbol\phi}_B}^2}{\,n-2p\,}.
\]
The same $\phi$-parameters appear because $Y_t=X_{-t}$ is an AR$(p)$ with
identical coefficients and an innovation $\tilde\nu_t\eqdist\eps_t$.
\end{definition}

\begin{definition}{Forward--backward least-squares estimator}{p4-fbls}
The \textbf{forward--backward} estimator $\hat{\boldsymbol\phi}_{FB}$ is the
minimiser of the \emph{combined} criterion
\[
\mathrm{SS}_F(\boldsymbol\phi)+\mathrm{SS}_B(\boldsymbol\phi).
\]
Equivalently it is the OLS solution of the stacked system obtained by piling
$\mathbf X_F,\mathbf X_B$ and $F,B$ on top of each other. Simulation studies
show it typically outperforms both the forward and the backward estimator.
\end{definition}

\begin{definition}{Least squares for general ARMA via residual reconstruction}{p4-lsarma}
For an ARMA$(p,q)$ model $\Phi(\Bsh)X_t=\Theta(\Bsh)\eps_t$ the innovations
$\{\eps_t\}$ are unobserved. The least-squares estimator chooses the parameters
minimising $S(\boldsymbol\phi,\boldsymbol\theta)=\sum_t\hat\eps_t^{\,2}$, where
for each candidate parameter value the residuals are \textbf{reconstructed
recursively} from the model equation after fixing the $q$ pre-sample
innovations (and any pre-sample data) to zero, e.g.\
$\eps_0=\eps_{-1}=\cdots=0$. The minimisation is nonlinear (no closed form)
except in the pure-AR case, where it reduces to the linear forward least squares
above.
\end{definition}

% ----------------------------------------------------------------------
\subsection{Theorems \& Propositions}

\begin{proposition}{Yule--Walker equations}{p4-yweq}
Let $\{X_t\}$ be a mean-zero causal AR$(p)$ process with ACVS $\acvs_\tau$.
Then for every $k\ge1$,
\[
\acvs_k=\sum_{j=1}^p\phi_j\,\acvs_{k-j},
\]
and in particular, taking $k=1,\dots,p$ and using $\acvs_{-\tau}=\acvs_\tau$,
\[
\boldsymbol\acvs=\Gamma\boldsymbol\phi,
\qquad
\boldsymbol\acvs=(\acvs_1,\dots,\acvs_p)\Tr,\quad
\Gamma=(\acvs_{i-j})_{i,j=1}^p .
\]
Moreover $\acvs_0=\sum_{j=1}^p\phi_j\acvs_j+\sigma_\eps^2$.\\[2pt]
\textit{Proof idea:} multiply $X_t=\sum_j\phi_jX_{t-j}+\eps_t$ by $X_{t-k}$ and
take expectations. Because the process is \emph{causal}, $X_{t-k}$ depends only
on $\eps_{t-k},\eps_{t-k-1},\dots$, so for $k>0$ we have $\E[\eps_tX_{t-k}]=0$;
with mean zero $\E[X_{t-j}X_{t-k}]=\acvs_{k-j}$, giving the recursion. For the
variance, multiply by $X_t$ instead: the cross term becomes
$\E[\eps_tX_t]=\E[\eps_t(\sum_j\phi_jX_{t-j}+\eps_t)]=\sigma_\eps^2$ (all
$\E[\eps_tX_{t-j}]=0$, $j\ge1$).\\
\textit{Why:} these $p+1$ equations express the $p+1$ unknowns
$(\phi_1,\dots,\phi_p,\sigma_\eps^2)$ purely in terms of $\acvs_0,\dots,\acvs_p$;
replacing the population ACVS by its sample version yields the Yule--Walker
estimators. It is the canonical method-of-moments fit for AR models.
\end{proposition}

\begin{proposition}{Minimising the prediction MSE gives Yule--Walker}{p4-mse}
Let $\{Y_t\}$ be mean-zero stationary with ACVS $\acvs$. The coefficients
$\phi_1,\dots,\phi_k$ minimising the one-step mean-square prediction error
\[
\E\Big[\big(Y_t-\textstyle\sum_{i=1}^k\phi_iY_{t-i}\big)^2\Big]
\]
satisfy the Yule--Walker equations $\acvs_j=\sum_{i=1}^k\phi_i\,\acvs_{j-i}$ for
$j=1,\dots,k$.\\[2pt]
\textit{Proof idea:} since $\{Y_t\}$ is mean-zero, the objective equals the
variance
$\acvs_0-2\sum_i\phi_i\acvs_i+\sum_{i,j}\phi_i\phi_j\acvs_{i-j}$, a positive
quadratic in $\boldsymbol\phi$. Setting $\partial/\partial\phi_i=0$ gives
$-2\acvs_i+2\sum_j\phi_j\acvs_{i-j}=0$, i.e.\ the Yule--Walker equations.\\
\textit{Why:} it shows Yule--Walker is not merely a moment trick but the
\emph{best linear predictor}: the AR coefficients are exactly the optimal
linear-prediction weights, which is why YW and least squares agree
asymptotically for AR models.
\end{proposition}

\begin{proposition}{Forward / backward least squares are the OLS estimators}{p4-ols}
With the design matrices of Definitions~\ref{def:p4-fls}--\ref{def:p4-bls}, the
minimisers of $\mathrm{SS}_F$ and $\mathrm{SS}_B$ are
\[
\hat{\boldsymbol\phi}_F=(F\Tr F)^{-1}F\Tr\mathbf X_F,
\qquad
\hat{\boldsymbol\phi}_B=(B\Tr B)^{-1}B\Tr\mathbf X_B,
\]
with residual-variance estimators
$\hat\sigma^2=\norm{\text{residuals}}^2/(n-2p)$.\\[2pt]
\textit{Proof idea:} each is a standard linear-regression sum of squares
$\norm{\mathbf y-A\boldsymbol\phi}^2$; the normal equations
$A\Tr A\,\boldsymbol\phi=A\Tr\mathbf y$ give the stated solution, and the
mean-square residual divides by the degrees of freedom $(n-p)-p$.\\
\textit{Why:} it places AR estimation inside the familiar OLS framework, so all
regression machinery (existence, uniqueness when $A\Tr A$ is invertible,
$\hat\sigma^2$) applies directly.
\end{proposition}

\begin{proposition}{Time reversibility of a Gaussian AR process}{p4-timerev}
If $\{X_t\}$ is a stationary Gaussian AR$(p)$ process with coefficients
$\phi_1,\dots,\phi_p$ and innovation variance $\sigma_\eps^2$, then the
time-reversed process $Y_t=X_{-t}$ is an AR$(p)$ process with the
\emph{same} coefficients and innovation variance:
\[
Y_t=\sum_{j=1}^p\phi_jY_{t-j}+\tilde\nu_t,\qquad \tilde\nu_t\eqdist\eps_t,
\]
equivalently $X_t=\sum_{j=1}^p\phi_jX_{t+j}+\nu_t$ with $\nu_t=\tilde\nu_{-t}$.\\[2pt]
\textit{Proof idea:} a stationary Gaussian process is determined by its ACVS,
and the ACVS is symmetric, $\acvs_{-\tau}=\acvs_\tau$; hence the forward and
reversed processes have identical second-order structure and (being Gaussian)
identical distributions. The full argument is given later via frequency-domain
methods (the spectral density is even in $f$).\\
\textit{Why:} it justifies the backward least-squares regression --- predicting
the present from the future uses the very same parameters --- and underpins the
forward--backward estimator.
\end{proposition}

\begin{proposition}{Levinson--Durbin: solving Yule--Walker in $O(p^2)$}{p4-ld}
The Yule--Walker system
$\hat\Gamma_p\hat{\boldsymbol\phi}_p=\hat{\boldsymbol\acvs}_p$ has a symmetric
Toeplitz coefficient matrix. The \textbf{Levinson--Durbin algorithm} exploits
this structure to solve it recursively in $k=1,\dots,p$, costing $O(p^2)$
operations rather than the $O(p^3)$ of a naive matrix inversion. As a
by-product it produces the partial autocorrelations
$\hat\pacf_k=\hat\phi_{k,k}$ and the prediction-error variances at every order.\\[2pt]
\textit{Proof idea:} the order-$k$ solution is obtained from the order-$(k-1)$
solution by one rank-one update governed by the reflection coefficient
$\hat\phi_{k,k}=(\hat\acvs_k-\sum_{j=1}^{k-1}\hat\phi_{k-1,j}\hat\acvs_{k-j})/
v_{k-1}$, where $v_{k-1}$ is the running error variance; the remaining
coefficients update as
$\hat\phi_{k,j}=\hat\phi_{k-1,j}-\hat\phi_{k,k}\hat\phi_{k-1,k-j}$.\\
\textit{Why:} for high-order AR fits, spectral estimation, or repeatedly
solving nested systems, the quadratic cost (and the free PACF) makes
Levinson--Durbin the standard solver.
\end{proposition}

\begin{proposition}{Least-squares fits need not be stationary}{p4-nonstat}
The forward, backward and forward--backward least-squares estimators
$\hat{\boldsymbol\phi}_F,\hat{\boldsymbol\phi}_B,\hat{\boldsymbol\phi}_{FB}$ are
\emph{not} constrained to yield a stationary (causal) model: the implied
polynomial $\hat\Phi(z)$ may have roots inside the unit circle.\\[2pt]
\textit{Proof idea:} the OLS objective places no constraint on the root moduli;
it only minimises the residual sum of squares, which can be smaller for a
non-causal parameter vector.\\
\textit{Why:} non-stationary estimates are a problem for \emph{prediction}, but
acceptable for some uses such as \emph{spectral estimation}, where they still
produce a valid spectral density estimate. The Yule--Walker estimator, by
contrast, always yields a stationary fit.
\end{proposition}

% ----------------------------------------------------------------------
% ----------------------------------------------------------------------
